\section{Completación graduada de los canales y principio de exclusión}
\label{vi:sec:fock-pauli}

La construcción de una partícula proporciona un espacio de canales, sus
etiquetas y los operadores que transportan orientación y memoria. La
composición de varias partículas añade una cuestión distinta: cómo se
intercambian esos canales y cómo se cuentan sus ocupaciones. Aquí se
construye ese estrato a partir de los datos ya producidos por
APP--TRIT--TPK. La masa y el Hamiltoniano pueden actuar después sobre
el espacio obtenido; sus valores no intervienen en la demostración de
las reglas de ocupación.

La paridad empleada es la del transporte de hoja en la representación
de una partícula. No es el signo de la carga, la paridad de un entero
visible ni la tercera coordenada del TRIT considerada aisladamente.
Las distintas lecturas conservan sus mapas. La
sección~\ref{vi:sec:conjugacion} construye el retorno orientado y la
representación espinorial; esta sección desarrolla el álgebra
multiparticular que corresponde a una graduación y una regla de
intercambio declaradas.

\subsection{Espacio finito de canales y carácter de paridad}
\label{vi:subsec:canales-graduados}

Fijada una profundidad \(K\), sea \(S_K\) el conjunto finito de
canales de la realización considerada. Sus elementos conservan los
datos de ruta y de representación que todavía distinguen estados:
hoja, orientación, frontera, memoria, sabor, color y componente de
espín cuando pertenecen al sector. El cociente que define \(S_K\)
se efectúa antes de construir la ocupación; no identifica dos canales
por la sola igualdad de sus masas.

El espacio de una partícula es
\begin{equation}
 \mathcal H_K=\bigoplus_{s\in S_K}V_s,
 \qquad\langle e_{s,a},e_{t,b}\rangle=\delta_{st}\delta_{ab},
 \label{vi:eq:espacio-canales-fock}
\end{equation}
donde \(V_s\) porta la multiplicidad de representación declarada.
Cada \(e_{s,a}\) es un modo completo, no una celda del atlas sin
etiquetas. Su producto interno positivo pertenece a esa realización
finita y hace explícita la normalización usada a continuación.

El dato de paridad es un carácter composicional del transporte,
\begin{equation}
 p:\operatorname{Mor}(\mathsf P_{\mathrm{TPK}})
 \longrightarrow\mathbb Z/2\mathbb Z,
 \qquad p(\gamma\eta)=p(\gamma)+p(\eta).
 \label{vi:eq:caracter-paridad-fock}
\end{equation}
Sobre los canales en los que está realizada esa graduación, la
involución de holonomía \(J_p\) satisface
\(J_pe=(-1)^{p(e)}e\). Por tanto,
\[
 \mathcal H_K=\mathcal H_{K,\bar0}\oplus\mathcal H_{K,\bar1},
 \qquad \Pi_{\bar0}=\frac{I+J_p}{2},\quad
 \Pi_{\bar1}=\frac{I-J_p}{2}.
\]
Los dos proyectores son ortogonales porque \(J_p^*=J_p\) y
\(J_p^2=I\). La igualdad \eqref{vi:eq:caracter-paridad-fock} asegura
que componer transportes conserva la graduación. Cuando una realización
no ha suministrado ese carácter, la completación que sigue es una
construcción relativa a su elección explícita, no una atribución
estadística obtenida de la apariencia de una ruta.

\subsection{Intercambio graduado y proyectores de permutación}
\label{vi:subsec:intercambio-graduado}

Para vectores homogéneos se define el intercambio
\begin{equation}
 \tau_p(v\otimes w)=(-1)^{p(v)p(w)}w\otimes v.
 \label{vi:eq:intercambio-graduado}
\end{equation}
El signo es negativo exactamente cuando ambos modos pertenecen al
sector impar. El intercambio no modifica las etiquetas de un modo:
transporta el modo completo de una posición tensorial a otra.

\Needspace{10\baselineskip}
\begin{proposition}[Representación graduada de las permutaciones]
\label{vi:prop:permutaciones-graduadas}
Los intercambios adyacentes \(\tau_{p,j}\) sobre
\(\mathcal H_K^{\otimes n}\) son unitarios y satisfacen
\[
\begin{gathered}
 \tau_{p,j}^2=I,\qquad
 \tau_{p,j}\tau_{p,k}=\tau_{p,k}\tau_{p,j}\quad(|j-k|>1),\\
 \tau_{p,j}\tau_{p,j+1}\tau_{p,j}
 =\tau_{p,j+1}\tau_{p,j}\tau_{p,j+1}.
\end{gathered}
\]
Definen, por tanto, una representación unitaria \(U_p\) de
\(\mathfrak S_n\).
\end{proposition}
\begin{proof}
Aplicar dos veces un intercambio da el signo
\((-1)^{2p(v)p(w)}=1\) y restaura ambos vectores. Los operadores
que actúan en pares disjuntos conmutan. Sobre una terna homogénea
de grados \(a,b,c\), los dos miembros de la relación de tres
intercambios producen el mismo orden final y el signo
\((-1)^{ab+ac+bc}\). Cada operador permuta una base ortonormal
multiplicándola por signos de módulo uno, luego es unitario.
\end{proof}

El promedio
\[
 P_n^p=\frac1{n!}\sum_{\pi\in\mathfrak S_n}U_p(\pi)
\]
es el proyector ortogonal sobre los tensores invariantes por el
intercambio graduado. En un espacio homogéneo par coincide con el
simetrizador ordinario; en uno homogéneo impar coincide con el
antisimetrizador. Si \(U\) denota la permutación sin signo,
\begin{equation}
 S_n=\frac1{n!}\sum_\pi U(\pi),\qquad
 A_n=\frac1{n!}\sum_\pi\operatorname{sgn}(\pi)U(\pi).
 \label{vi:eq:simetrizadores-fock}
\end{equation}

\begin{proposition}[Simetrización y antisimetrización exactas]
\label{vi:prop:proyectores-fock}
Se tiene \(S_n^*=S_n^2=S_n\) y \(A_n^*=A_n^2=A_n\).
Para \(n\ge2\), \(S_nA_n=0\). Sus imágenes son, respectivamente,
\(\operatorname{Sym}^n\mathcal H\) y \(\bigwedge^n\mathcal H\).
\end{proposition}
\begin{proof}
La adjunción sustituye cada permutación por su inversa y deja
invariantes las sumas. En el cuadrado de cada promedio, toda
permutación \(\rho\) aparece \(n!\) veces; en el segundo promedio
su coeficiente es \(\operatorname{sgn}(\rho)\). Esto prueba
la idempotencia. El producto de ambos proyectores incluye la suma
de los signos de todas las permutaciones, que es cero para
\(n\ge2\). Finalmente, sus imágenes son los tensores invariantes
y alternantes por sus propias fórmulas de promedio.
\end{proof}

\subsection{Álgebra universal y normalización de los estados de ocupación}
\label{vi:subsec:fock-construccion}

Sea \(T(\mathcal H)=\bigoplus_{n\ge0}\mathcal H^{\otimes n}\)
el álgebra tensorial. La completación algebraica graduada es el
cociente por el ideal generado por
\begin{equation}
 v\otimes w-(-1)^{p(v)p(w)}w\otimes v.
 \label{vi:eq:ideal-fock}
\end{equation}
La propiedad universal del cociente significa que toda aplicación
lineal graduada de los canales a un álgebra que cumpla estas
relaciones se prolonga de manera única a un homomorfismo del
cociente. La regla de intercambio, una vez declarada, fija así el
producto de ocupaciones sin recurrir a energías objetivo.

\begin{theorem}[Descomposición simétrica y exterior]
\label{vi:thm:fock-graduado}
Si \(d_0=\dim\mathcal H_{\bar0}\) y
\(d_1=\dim\mathcal H_{\bar1}\), la completación graduada tiene
la descomposición
\begin{equation}
 \mathcal F(\mathcal H)=
 \mathcal F_+(\mathcal H_{\bar0})\otimes
 \mathcal F_-(\mathcal H_{\bar1}),
 \label{vi:eq:fock-producto}
\end{equation}
\[
 \mathcal F_+(\mathcal H_{\bar0})=
 \bigoplus_{m\ge0}\operatorname{Sym}^m\mathcal H_{\bar0},
 \qquad
 \mathcal F_-(\mathcal H_{\bar1})=
 \bigoplus_{k=0}^{d_1}\bigwedge^k\mathcal H_{\bar1}.
\]
El sector de número total \(n\) es
\[
 \mathcal F_n(\mathcal H)=
 \bigoplus_{k=0}^{\min(n,d_1)}
 \operatorname{Sym}^{n-k}\mathcal H_{\bar0}
 \otimes\bigwedge^k\mathcal H_{\bar1}.
\]
\end{theorem}
\begin{proof}
Fijadas bases homogéneas ordenadas, las relaciones
\eqref{vi:eq:ideal-fock} permiten colocar primero los factores pares
y después los impares. Los pares conmutan; los impares anticonmutan,
y el cuadrado de un impar es cero porque se trabaja en característica
cero. Cada monomio queda escrito como
\[
 e_1^{m_1}\cdots e_{d_0}^{m_{d_0}}
 f_{i_1}\cdots f_{i_k},
 \qquad m_j\ge0,\quad i_1<\cdots<i_k.
\]
La realización mediante tensores simétricos y alternantes hace
linealmente independientes esos monomios. Proporciona la aplicación
inversa y demuestra la descomposición. Agruparlos por
\(m_1+\cdots+m_{d_0}+k=n\) da el último enunciado.
\end{proof}

La completación de Hilbert declara ortonormales los vectores de
ocupación normalizados \(|\boldsymbol m;I\rangle\), donde
\(\boldsymbol m\in\mathbb N_0^{d_0}\) e
\(I\subset\{1,\ldots,d_1\}\). En el sector bosónico de grado
\(n=\sum_jm_j\), la normalización tensorial es
\[
 |\boldsymbol m\rangle=
 \sqrt{\frac{n!}{\prod_jm_j!}}\,
 S_n(e_1^{\otimes m_1}\otimes\cdots\otimes
 e_{d_0}^{\otimes m_{d_0}}).
\]
Para \(I=\{i_1<\cdots<i_k\}\), el vector exterior es
\[
 |I\rangle=\frac1{\sqrt{k!}}\sum_{\pi\in\mathfrak S_k}
 \operatorname{sgn}(\pi)
 f_{i_{\pi(1)}}\otimes\cdots\otimes f_{i_{\pi(k)}}.
\]
En la primera fórmula hay \(n!/\prod m_j!\) tensores distintos;
en la segunda hay \(k!\). Contarlos prueba que ambas normas son uno.
El grado cero es \(\mathbb C|0\rangle\); este vacío de ocupación
no se identifica con la carta constitutiva del vacío electromagnético
ni con una vacancia del registro TPK.

\subsection{Creación fermiónica y derivación de las relaciones CAR}
\label{vi:subsec:car-construccion}

Sea \(I\) un subconjunto ordenado de modos impares y
\(r_j(I)=\#\{i\in I:i<j\}\). Se definen
\begin{align}
 a_j^\dagger|I\rangle&=
 \begin{cases}
 0,&j\in I,\\
 (-1)^{r_j(I)}|I\cup\{j\}\rangle,&j\notin I,
 \end{cases}\notag\\
 a_j|I\rangle&=
 \begin{cases}
 (-1)^{r_j(I)}|I\setminus\{j\}\rangle,&j\in I,\\
 0,&j\notin I.
 \end{cases}
 \label{vi:eq:creacion-fermi}
\end{align}
El signo cuenta las transposiciones necesarias para insertar o retirar
el modo manteniendo el orden de la base exterior. No usa su masa,
su carga ni el signo de un valor escalar publicado. Las matrices de
ambas aplicaciones son adjuntas y sus normas son \(1\), salvo el
caso trivial sin modos impares.

\begin{theorem}[Relaciones canónicas de anticonmutación]
\label{vi:thm:car-construidas}
Los operadores de \eqref{vi:eq:creacion-fermi} satisfacen
\[
 \{a_i,a_j^\dagger\}=\delta_{ij}I,\qquad
 \{a_i,a_j\}=0,\qquad
 \{a_i^\dagger,a_j^\dagger\}=0.
\]
Son relaciones exactas sobre la totalidad de
\(\mathcal F_-(\mathcal H_{\bar1})\).
\end{theorem}
\begin{proof}
Para \(i=j\), el producto \(a_j^\dagger a_j\) es la identidad
en los vectores que contienen \(j\) y cero en los restantes;
\(a_ja_j^\dagger\) tiene las acciones complementarias. Su suma
es \(I\). Dos inserciones o dos retiradas del mismo modo dan cero.
Para \(i\ne j\), si alguna operación está impedida, ambos términos
del anticomutador correspondiente son cero. Si ambas composiciones
son admisibles, producen el mismo conjunto de ocupaciones. Cambiar
el orden modifica en una unidad el número de modos que preceden a
una de las inserciones o retiradas; los signos son opuestos. Esto
prueba los tres anticomutadores en cada vector de la base.
\end{proof}

\begin{corollary}[Exclusión y ocupación de un modo completo]
\label{vi:cor:pauli-ocupacion}
Para \(N_j=a_j^\dagger a_j\),
\[
 N_j^*=N_j,\qquad N_j^2=N_j,\qquad
 \operatorname{Spec}(N_j)\subset\{0,1\}.
\]
El mismo modo completo admite como máximo una ocupación fermiónica.
\end{corollary}
\begin{proof}
Las CAR dan
\[
 N_j^2=a_j^\dagger a_ja_j^\dagger a_j
 =a_j^\dagger(I-a_j^\dagger a_j)a_j=N_j,
\]
porque \(a_j^2=(a_j^\dagger)^2=0\). La autoadjunción es inmediata;
un autovalor de un idempotente autoadjunto vale \(0\) o \(1\).
\end{proof}

La prueba construye CAR antes de deducir Pauli: no introduce la
exclusión como restricción adicional de una tabla de partículas.
También explica los generadores reales empleados en la
proposición~\ref{vi:prop:car-real}: esa proposición actúa ahora
sobre el álgebra exterior explícita. Su interpretación Majorana
conserva las condiciones de la sección~\ref{vi:sec:topologia}.

\subsection{Creación bosónica, CCR y dominio de partículas finitas}
\label{vi:subsec:ccr-construccion}

Sobre la base bosónica normalizada se definen
\begin{equation}
 b_j^\dagger|\boldsymbol m\rangle=
 \sqrt{m_j+1}\,|\boldsymbol m+\boldsymbol e_j\rangle,
 \qquad
 b_j|\boldsymbol m\rangle=
 \sqrt{m_j}\,|\boldsymbol m-\boldsymbol e_j\rangle,
 \label{vi:eq:creacion-bose}
\end{equation}
con valor cero en la segunda expresión cuando \(m_j=0\).
Los factores proceden de la normalización de tensores simétricos.
El dominio común \(\mathcal D_{\mathrm{fin}}\) es la suma directa
algebraica de los sectores de número: los vectores tienen sólo un
número finito de componentes de ocupación no nulas. Es denso e
invariante por las aplicaciones anteriores.

\begin{theorem}[Relaciones canónicas de conmutación]
\label{vi:thm:ccr-construidas}
En \(\mathcal D_{\mathrm{fin}}\),
\[
 [b_i,b_j^\dagger]=\delta_{ij}I,\qquad
 [b_i,b_j]=[b_i^\dagger,b_j^\dagger]=0.
\]
Los operadores de número \(N_j^B=b_j^\dagger b_j\) satisfacen
\(N_j^B|\boldsymbol m\rangle=m_j|\boldsymbol m\rangle\).
\end{theorem}
\begin{proof}
Para un mismo modo,
\(b_jb_j^\dagger|\boldsymbol m\rangle=(m_j+1)|\boldsymbol m\rangle\)
y \(b_j^\dagger b_j|\boldsymbol m\rangle=m_j|\boldsymbol m\rangle\).
La diferencia es la identidad. Para modos distintos, los factores
afectan coordenadas de ocupación diferentes, por lo que conmutan.
Las fórmulas de número siguen de aplicar consecutivamente las dos
operaciones de \eqref{vi:eq:creacion-bose}.
\end{proof}

Si hay algún modo par, \(b_j\) y \(b_j^\dagger\) son operadores
no acotados en la completación de Hilbert: sus normas sobre
\(|m_j\rangle\) crecen como raíces de \(m_j\). Las igualdades
anteriores tienen por ello un dominio explícito y no se presentan
como identidades entre matrices finitas en todo el espacio bosónico.
Los operadores fermiónicos actúan sobre el segundo factor de
\eqref{vi:eq:fock-producto}; los bosónicos, sobre el primero. En
el dominio común, los operadores de esos dos factores conmutan.

\begin{proposition}[Defecto exacto de un corte bosónico duro]
\label{vi:prop:corte-bosonico}
En un modo, sea \(P_N\) la proyección sobre
\(|0\rangle,\ldots,|N\rangle\) y sea \(b^{(N)}=P_NbP_N\).
Sobre ese espacio finito,
\[
 [b^{(N)},(b^{(N)})^\dagger]
 =I-(N+1)|N\rangle\langle N|.
\]
\end{proposition}
\begin{proof}
Para \(m<N\), la creación y aniquilación conservan el cálculo
\((m+1)-m=1\). En \(|N\rangle\), la creación truncada es
cero y el producto inverso vale \(N\); el conmutador vale
\(-N=1-(N+1)\). Son todos los vectores de la base.
\end{proof}

La corrección reside en el último nivel del corte, no en la ley de
ocupación. Un cálculo finito puede verificar las CCR lejos de esa
frontera y debe conservar el término de borde cuando incluye el
nivel terminal. Esta distinción es indispensable para que un
certificado matricial no confunda un truncamiento con la
completación bosónica completa.

\subsection{Qué significa «el mismo estado» en el principio de Pauli}
\label{vi:subsec:pauli-etiquetas-hojas}

La exclusión se refiere al mismo vector de un modo completo. Dos
canales pueden compartir celda, masa o firma reducida y seguir
siendo linealmente independientes por su espín, hoja, color,
sabor u otra coordenada que el transporte conserve. Esos dos
vectores admiten un producto exterior no nulo. Si un cociente
físico los identifica, ese cociente debe actuar antes de afirmar
cuántos modos independientes hay.

\begin{proposition}[Criterio de independencia para una ocupación doble]
\label{vi:prop:pauli-gram}
Para dos vectores de un espacio de canales,
\[
 \|u\wedge v\|^2
 =\|u\|^2\|v\|^2-|\langle u,v\rangle|^2.
\]
En particular, dos modos ortonormales producen un estado exterior
de norma uno y dos vectores proporcionales producen cero.
\end{proposition}
\begin{proof}
La realización exterior de grado dos es
\((u\otimes v-v\otimes u)/\sqrt2\). Expandir su norma da
la fórmula. La igualdad en Cauchy--Schwarz caracteriza la
dependencia lineal; la ortonormalidad da norma cuadrada uno.
\end{proof}

Para un canal orbital \(u\) y dos componentes de espín
ortogonales \(\xi_+,\xi_-\), los modos
\(u\otimes\xi_+\) y \(u\otimes\xi_-\) son distintos y su
producto exterior puede estar ocupado. No existe doble ocupación
de un mismo modo: existe ocupación de dos componentes distintas
de la representación. Lo mismo se aplica a dos hojas cuando éstas
siguen siendo componentes independientes de la realización.

La palabra electrónica central puede ser fija por \(C_M\), mientras
la representación de carga distingue electrón y positrón. Esa
coincidencia de palabra no identifica sus modos completos ni
decide su ocupación. La involución de palabra, la inversión
celular y el intercambio de factores de Fock conservan los
dominios diferenciados de la sección~\ref{vi:sec:conjugacion}.

Si \(r\) modos fermiónicos ortonormales poseen el mismo valor
de un observable parcial, su ocupación total puede tomar valores
\(0,1,\ldots,r\). Cada \(N_j\) sigue siendo un proyector;
la suma \(\sum_{j=1}^rN_j\) no tiene por qué serlo. Esta
consecuencia distingue degeneración espectral y exclusión de
un modo, dos propiedades que una cifra de masa aislada no separa.

\subsection{Dimensiones finitas y transporte por refinamiento}
\label{vi:subsec:fock-refinamiento}

\begin{proposition}[Censo de ocupaciones]
\label{vi:prop:censo-fock}
Para \(d\) modos,
\[
 \dim\bigwedge^n\mathbb C^d=\binom dn,
 \qquad
 \dim\operatorname{Sym}^n\mathbb C^d=\binom{d+n-1}{n}
 \quad(d\ge1).
\]
La primera dimensión es cero para \(n>d\). En el espacio
graduado, las dimensiones de número total \(n\) se obtienen
sumando los productos de ambas dimensiones para
\(k=0,\ldots,\min(n,d_1)\).
\end{proposition}
\begin{proof}
Una base exterior se elige mediante subconjuntos de \(n\) modos
distintos. Una base simétrica se elige mediante
\((m_1,\ldots,m_d)\in\mathbb N_0^d\) de suma \(n\).
La colocación de \(d-1\) separadores entre \(n\) ocupaciones
cuenta esas composiciones débiles y da el segundo binomio.
La fórmula graduada sigue de \eqref{vi:eq:fock-producto} por
sumas directas y productos tensoriales.
\end{proof}

Sean \(T:\mathcal H_K\to\mathcal H_L\) los entrelazadores
de refinamiento declarados, que conservan la graduación y los
datos de canal correspondientes. Cada uno induce
\begin{equation}
 \Gamma(T)=\bigoplus_{n\ge0}T^{\otimes n}|_{\mathcal F_n},
 \qquad T^{\otimes0}=I_{\mathbb C}.
 \label{vi:eq:segunda-cuantizacion-transporte}
\end{equation}

\begin{theorem}[Compatibilidad de la completación con el transporte]
\label{vi:thm:fock-naturalidad}
Si \(T\) conserva el grado, entonces
\(T^{\otimes n}P_n^p=P_n^{p'}T^{\otimes n}\).
Si además \(\|T\|\le1\), \(\Gamma(T)\) es un operador
acotado de norma a lo sumo \(1\) sobre Fock y
\[
 \Gamma(T_2T_1)=\Gamma(T_2)\Gamma(T_1),\qquad
 \Gamma(I)=I.
\]
Por consiguiente, los cuadrados compatibles de refinamiento
de los canales inducen cuadrados compatibles de ocupaciones.
\end{theorem}
\begin{proof}
La conservación del grado hace conmutar \(T^{\otimes n}\)
con cada intercambio adyacente; por la
proposición~\ref{vi:prop:permutaciones-graduadas}, conmuta
con las permutaciones y con su promedio. Las restricciones
quedan así bien definidas. Su norma está acotada por
\(\|T\|^n\le1\), de modo que la suma directa es acotada.
Las identidades y composiciones se verifican en cada grado
y, por continuidad, en la completación.
\end{proof}

La contractividad es una condición suficiente que permite tratar
la suma bosónica sin restricciones adicionales. En grados finitos
las potencias existen para cualquier aplicación lineal, pero su
suma no siempre es acotada: para un modo bosónico y \(T=2I\),
\(\Gamma(T)|n\rangle=2^n|n\rangle\). En el sector exterior
finito sólo hay \(d_1+1\) grados, de modo que cualquier
aplicación lineal de ese sector tiene una segunda cuantización
finita acotada. La diferencia es de dominio, no de causalidad.

\begin{proposition}[Transporte de creaciones y aniquilaciones]
\label{vi:prop:fock-transporte-adjunto}
Sea \(T:\mathcal H_K\to\mathcal H_L\) un transporte acotado
que conserva el grado. Para \(f\in\mathcal H_K\) y
\(g\in\mathcal H_L\), en el dominio algebraico de partículas
finitas se cumplen
\begin{align}
 \Gamma(T)c_K^\dagger(f)
 &=c_L^\dagger(Tf)\Gamma(T),
 \label{vi:eq:fock-transporte-creacion}\\
 c_L(g)\Gamma(T)
 &=\Gamma(T)c_K(T^*g).
 \label{vi:eq:fock-transporte-aniquilacion}
\end{align}
Los subíndices indican la fibra sobre la que actúa cada operador.
Si \(T\) es isométrico, la segunda identidad se especializa a
\[
 c_L(Tf)\Gamma(T)=\Gamma(T)c_K(f),
 \qquad
 N_{L,Tf}\Gamma(T)=\Gamma(T)N_{K,f},
\]
donde \(N_{K,f}=c_K^\dagger(f)c_K(f)\) y, para un modo unitario,
\(N_{L,Tf}\) es su operador de ocupación transportado.
\end{proposition}
\begin{proof}
En un producto simétrico o exterior, aplicar \(T\) a cada factor
conmuta con insertar \(f\), que se convierte en \(Tf\). Esto
demuestra \eqref{vi:eq:fock-transporte-creacion}. La aniquilación
contrae cada factor con el modo observado. Para cada factor \(v\),
\[
 \langle g,Tv\rangle=\langle T^*g,v\rangle.
\]
Los signos exteriores, o graduados, se conservan porque \(T\)
preserva el grado; las normalizaciones simétricas son iguales en
ambos miembros. La suma de contracciones demuestra
\eqref{vi:eq:fock-transporte-aniquilacion}. Si \(T^*T=I\), se
sustituye \(g=Tf\). La composición de las identidades de creación
y aniquilación da después la de ocupación. Todas las operaciones
consideradas preservan el dominio de partículas finitas; cuando
\(\|T\|\le1\), \(\Gamma(T)\) posee además la extensión acotada
del teorema~\ref{vi:thm:fock-naturalidad}.
\end{proof}

La identidad general conserva el adjunto \(T^*\) y es válida sin
isometría. La necesidad de la condición isométrica aparece al
transportar el mismo modo. Por ejemplo, si \(T=tI\), con
\(0<t<1\), para un modo unitario \(f\) y el vacío \(\Omega\),
\[
 c(tf)\Gamma(T)f=t^2\Omega,
 \qquad \Gamma(T)c(f)f=\Omega.
\]
La fórmula general da correctamente
\(\Gamma(T)c(T^*tf)f=t^2\Omega\). Se conserva así el
transporte por contracciones, con su adjunto explícito, y la
especialización isométrica de la ocupación.

\subsection{Simetría de intercambio y registro de un compuesto}
\label{vi:subsec:pauli-composicion}

La composición de la sección~\ref{vi:subsec:composicion-registro}
actúa primero sobre historias y sus fronteras. Fock organiza
la ocupación de canales; \(\mathfrak C_{12}\) organiza el
registro de enlace. Son operaciones coordinadas que no se
sustituyen entre sí. En particular, la antisimetrización no
suprime la memoria del orden de acoplamiento.

Sea \(\mathcal D\subset\mathcal H^{\otimes n}\) el espacio
generado por expresiones de constituyentes admisibles para una
interfaz. La linealización de la composición se escribe
\(C_{\mathcal T}:\mathcal D\to\mathcal H_{\mathrm{comp}}\).
Cuando un grupo finito \(G\) de intercambios conserva la
admisibilidad y la frontera, la compatibilidad concreta es
\[
 C_{\mathcal T}U(g)=V(g)C_{\mathcal T},\qquad g\in G,
\]
donde \(U,V\) son representaciones unitarias y \(V\) transporta
los registros compuestos, incluidos sus datos de hoja y memoria.
La ecuación es la condición de
intercambio de esa realización; no se infiere del simple
censo de expresiones.

\begin{proposition}[Descenso de la composición al sector estadístico]
\label{vi:prop:composicion-estadistica}
Sea \(\chi:G\to\{+1,-1\}\) el carácter del sector de
intercambio. Bajo la compatibilidad anterior,
\[
 P_\chi^{\mathrm{comp}}C_{\mathcal T}
 =C_{\mathcal T}P_\chi^{\mathrm{in}},
\]
donde
\[
 P_\chi^{\mathrm{in}}=\frac1{|G|}\sum_g\chi(g)U(g),
 \qquad
 P_\chi^{\mathrm{comp}}=\frac1{|G|}\sum_g\chi(g)V(g).
\]
\end{proposition}
\begin{proof}
Multiplicar cada igualdad de entrelazamiento por
\(\chi(g)/|G|\) y sumar da exactamente la identidad.
Los promedios son proyectores por el mismo cálculo de
la proposición~\ref{vi:prop:proyectores-fock}.
\end{proof}

El resultado justifica componer dentro de un sector simétrico
o alternante cuando la frontera posee el transporte requerido.
Si una permutación cambia el canal físico de enlace, se
conservan los árboles distintos y sus interrelaciones; no se
identifican mediante una simetrización que olvide esa frontera.

\subsubsection{Antisimetría cromática y resto del estado bariónico}

En la realización bariónica, la
proposición~\ref{vi:prop:singlete-barion} construye el vector
cromático \(|\varepsilon\rangle\), que cambia de signo
bajo una transposición. Tras reagrupar los factores del estado
completo, un vector factorizado tiene la forma
\[
 \Psi=|\varepsilon\rangle_{\mathrm{col}}
 \otimes\Phi_{\mathrm{espín,sabor,orbital,hoja}}.
\]
El intercambio de dos constituyentes actúa simultáneamente
en ambos factores. Si \(U_{\mathrm{rest}}(\pi)\Phi=\Phi\),
entonces el estado completo cambia de signo y pertenece al
sector fermiónico. Recíprocamente, para un vector factorizado
con el singlete cromático no nulo, la alternancia total exige
que el resto sea simétrico bajo esas mismas permutaciones.
La prueba consiste en factorizar
\(U_{\mathrm{tot}}(\pi)=U_{\mathrm{col}}(\pi)\otimes
U_{\mathrm{rest}}(\pi)\) y cancelar el signo cromático.

La palabra «resto» incluye explícitamente sabor y hoja.
Intercambiar colores sin transportar esas componentes no es
la permutación de dos modos completos. Para un estado no
factorizable se aplica el proyector alternante al espacio
tensorial completo, sin sustituirlo por una regla escalar.

\subsubsection{Compatibilidad de proyectores físicos}

Los proyectores de color, canal y estadística pueden combinarse
por producto cuando conmutan. Si \(P,Q\) son proyectores
ortogonales y \([P,Q]=0\), entonces
\((PQ)^*=PQ\), \((PQ)^2=PQ\) y
\(\operatorname{im}(PQ)=\operatorname{im}P\cap\operatorname{im}Q\).
La inclusión en ambas imágenes es directa; para un vector
en la intersección, \(PQv=v\). Si no conmutan, su producto
no representa en general el proyector de las dos condiciones
simultáneas: la realización debe construir el subespacio
común o el orden físico de las operaciones.

Esta precisión conserva el sentido de los proyectores
singlete y de Pauli que intervienen en los compuestos de la
sección~\ref{vi:sec:familias-compuestos}. El acoplamiento
selecciona un sector; la masa compuesta se evalúa después
de construir el registro de enlace de ese sector.

\subsection{Acción del operador material y alcance de la identificación física}
\label{vi:subsec:fock-alcance}

Sea \(B=B^*\) un operador de un canal ya obtenido de la
dinámica interna y que conserva la graduación. En grado \(n\)
actúa mediante
\[
 d\Gamma_n(B)=\sum_{j=1}^n
 I^{\otimes(j-1)}\otimes B\otimes I^{\otimes(n-j)}.
\]
La suma conmuta con las permutaciones graduadas y se restringe
a \(\mathcal F_n\). Si \(B\ge0\), cada restricción es positiva.
En una base propia, su valor es la suma de las ocupaciones
por los valores propios de \(B\), con las ocupaciones
permitidas por las pruebas anteriores. Esta es la acción
de un observable aditivo sobre canales independientes;
no identifica el operador de un compuesto ligado con la
suma de las masas aisladas de sus constituyentes.

En la suma de Hilbert bosónica, \(d\Gamma(B)\) posee el
dominio natural
\[
 \left\{\psi=(\psi_n):
 \sum_n\|d\Gamma_n(B)\psi_n\|^2<\infty\right\}.
\]
Cada bloque es autoadjunto en dimensión finita, y su suma
directa con ese dominio es autoadjunta. El operador de enlace
o de canales abiertos se añade sólo cuando ha sido construido
por la dinámica correspondiente. La formación de un polo
complejo sigue el tratamiento de
\ref{vi:subsec:resonancias-familias}.

Se han obtenido, por tanto, resultados algebraicos completos:
proyectores de simetría, bases normalizadas, CAR y CCR en
sus dominios, exclusión de un modo, censos de ocupación y
compatibilidad con entrelazadores de refinamiento. El punto
de partida de esos resultados está declarado en
\eqref{vi:eq:espacio-canales-fock},
\eqref{vi:eq:caracter-paridad-fock} y
\eqref{vi:eq:intercambio-graduado}.

Cuando la realización de rotaciones satisface
\(U(2\pi)=(-1)^p\), la misma graduación coordina la
restauración de orientación a \(4\pi\) y la completación
exterior. Esa compatibilidad no se sustituye por una
coincidencia entre dos signos. La identificación con el
teorema relativista de espín--estadística requiere además
la representación lorentziana, la localidad y la positividad
de la realización física. Los sectores de trenzado que no
descienden a permutaciones permanecen en el codominio de
la sección~\ref{vi:sec:topologia}, no se fuerzan dentro
de esta completación bosónica o fermiónica.
